\section{Background and Research Questions}
\label{sec:background}

\subsection{Offline RCA Benchmarks}
\label{sec:bg-bench}

In microservice root cause analysis, a method receives telemetry from a window around a failure (metrics and, in some benchmarks, logs and traces) and returns a ranked list of candidate root causes, usually services or components.
A case is scored by whether the true root cause is ranked first (\acc) or, in \openrca, by the fraction of required answer items that are correct.
Public RCA evaluation happens in four modes: (a) offline benchmarks that release fault cases, ground truth, and a scoring rule, so that every method can be scored on the same cases (\openrca~\citep{xu2025openrca}, \rcaeval~\citep{pham2024rcaeval}, \petshop~\citep{hardt2024petshop}); (b) single-application testbeds such as Train Ticket~\citep{zhou2018trainticket}, which each paper deploys and injects anew; (c) competitions with year-specific protocols; and (d) interactive benchmarks in which an agent queries a live system or a recorded state snapshot through diagnostic tools~\citep{aiopslab,jha2025itbench,wang2026cloudopsbench}.
We audit mode (a), released static-telemetry datasets and their evaluation pipelines; interactive benchmarks are outside the scope of this audit.

An offline benchmark is released as per-case data tables, a runner that loads them and executes the methods, and a scoring script.
A pooled score is the product of that pipeline: which table the runner reads, what happens when a method fails on a case, which cases enter the denominator, and how cases are scored and weighted.
The benchmark papers state the scoring rule; the other choices are made in code.
Each benchmark also contains several applications or traffic scenarios, which we call subsystems; its paper and the method papers that use it report a score per subsystem and a score pooled over all of them, and it is the pooled score that is ranked.

\subsection{Research Questions}
\label{sec:rqs}

A pooled leaderboard reports one number per method. Reading it as evidence that method $A$ should be preferred to method $B$ assumes two things the number does not state: on which systems the advantage holds, and how the score was produced. We ask one question about each.

\begin{description}
\item[RQ1 (scope).] How stable are comparisons between published methods across the subsystems of a benchmark, and across a second, independently collected release of the same applications?
\item[RQ2 (provenance).] Which choices in the released evaluation pipeline---the input representation, the handling of failed outputs, the denominator, and the scoring and weighting rule---change these comparisons, and by how much?
\end{description}
